%%%%%%%%%%%%%%%%%%%%%%%%%%%%%%%%%%%%%%%
% Trabalho de Sistemas Distribu�dos
%%%%%%%%%%%%%%%%%%%%%%%%%%%%%%%%%%%%%%%

\documentclass[12pt,brazil, a4paper]{article}

%%%%%%%%%%%%%%%%%%%%%%%%%%%%%%%%%%%%%%%
% Preambulo - Configura��o de p�gina
%%%%%%%%%%%%%%%%%%%%%%%%%%%%%%%%%%%%%%%
\usepackage{color}
\usepackage[portuges]{babel}
\usepackage[T1]{fontenc}
\usepackage[ansinew]{inputenc}
\usepackage{array}
\usepackage{colortbl}
\usepackage{epsfig}
\usepackage{multirow}
\usepackage{hyperref}
\usepackage{amsmath}

\usepackage{tikz}
\usetikzlibrary{3d}
\usetikzlibrary{calc,decorations.pathmorphing} % noisy shapes
\usetikzlibrary{fit}					% fitting shapes to coordinates
\usetikzlibrary{backgrounds}	% drawing the background after the foreground
\usetikzlibrary{shadows}
\usetikzlibrary{mindmap}
\usetikzlibrary{positioning,shapes,shadows,arrows}
\usepackage{pgfplots}
\pgfplotsset{compat=1.8}
\usetikzlibrary{arrows,shapes}

\pagestyle{empty}

% Define margem no topo como sendo 1,7 cm
\setlength{\topmargin}{0pt}
\setlength{\voffset}{-1.8cm}

%define margem inferior como 1,5 cm
\setlength{\textheight}{25,6cm}

%define margens laterais como sendo 1,5 cada
\setlength{\oddsidemargin}{0cm}
\setlength{\hoffset}{-1.2cm}
\setlength{\textwidth}{18cm}

%Define campo com sombreamento cinza
\newcolumntype{G}{>{\columncolor[gray]{0.8}[\tabcolsep]\tiny \bf \sf}p{13pt}<{\centering}}
\newcolumntype{T}{>{\tiny \sf}p{13pt}<{\centering}}


\usepackage{xcolor}
\definecolor{verde}{rgb}{0,0.5,0}

\usepackage{listings}
\lstset{
	language=c,
	basicstyle=\ttfamily\small, 
	keywordstyle=\color{blue}, 
	stringstyle=\color{verde}, 
	commentstyle=\color{red}, 
	extendedchars=true, 
	showspaces=false, 
	showstringspaces=false, 
	numbers=left,
	numberstyle=\tiny,
	breaklines=true, 
	backgroundcolor=\color{green!10},
	breakautoindent=true, 
	captionpos=b,
	xleftmargin=0pt,
}


%%%%%%%%%%%%%%%%%%%%%%%%%%%%%%%%%%%%%%%
% Inicio do documento
%%%%%%%%%%%%%%%%%%%%%%%%%%%%%%%%%%%%%%%
\begin{document}
	
	\selectlanguage{brazil}
	
	\hspace{-0.8cm} \begin{tabular*}{18cm}{p{2cm}p{12cm}@{\extracolsep{\fill}}}
		\multirow{2}{2cm}{\epsfig{file=puclogo_small_bw,width=1.5cm}}
		&  \vspace{1mm} \usefont{T1}{cmss}{bx}{n} {\centering \footnotesize PONTIF�CIA UNIVERSIDADE CAT�LICA DE MINAS GERAIS }\\
		&  \usefont{T1}{cmss}{x}{n} Instituto de Ci�ncias Exatas e Inform�tica \vspace{3mm}  \\
	\end{tabular*}
	
	
	%%%%%%%%%%%%%%%%%%%%%%%%%%%%%%%%%%%%%%%
	% Define caracteristicas da disciplina
	%%%%%%%%%%%%%%%%%%%%%%%%%%%%%%%%%%%%%%%
	\hspace{-0.8cm} \begin{tabular*}{18.5cm}{|p{7cm}@{\extracolsep{\fill}}|p{6cm} p{1.7cm}|}
		\hline
		\sf Disciplina & \sf Professor  &  \\
		
		\multicolumn{1}{|r|}{Algoritmos e Estruturas de Dados} &
		\multicolumn{2}{c|}{Kleber Jacques F. de Souza (klebersouza@pucminas.br)} \\ \hline
		
	\end{tabular*}
	
	
	\vspace{0.1cm}
	
	
	\begin{center}
		{\Large \bf Avalia��o 1A} \\
		%{\normalsize \bf em dupla}
	\end{center}
	

\begin{flushleft}
	\textbf{Nome:}  \makebox[5.0in]{\hrulefill}\textbf{Nota:}  \makebox[0.85in]{\hrulefill}\\ 
\end{flushleft} 

\section*{Instru��es}

\begin{itemize}
\footnotesize
\item A interpreta��o das quest�es faz parte da avalia��o.
\item Coloque as observa��es que julgar necess�rias na solu��o das quest�es.
\end{itemize}

%===================================================================================================
\section*{\centering Fila de Espera com Prioridade}

Desenvolva um programa que seja capaz de gerenciar uma Fila de espera. O programa deve possuir as seguintes funcionalidades:

\begin{itemize}
	\item \textbf{Retirar uma senha}: o programa deve receber a idade e nome de uma pessoa e retornar uma senha de espera.
	\item \textbf{Chamar a pr�xima Pessoa}: o programa deve exibir a senha, o nome e idade da pr�xima pessoas da Fila a ser atendida.
	\item \textbf{Listar Fila de espera}: o programa deve exibir a lista com o nome e idade de todas as pessoas que est�o aguardando na Fila.
	\item \textbf{Listar tamanho da Fila}: o programa deve exibir a quantidade de pessoas esperando na Fila.
	\item \textbf{Gerenciar prioridades}: o programa deve garantir que pessoas acima de 80 anos ou mais tenham prioridades sobre pessoas mais novas que elas, e que pessoas de 65 anos ou mais tenham prioridades sobre pessoas mais novas que elas.
\end{itemize}

Voc� pode desenvolver quaisquer outras funcionalidades que julgar necess�rias.



\begin{flushright}
	\textbf{Boa Prova!} 
\end{flushright} 



%===================================================================================================

\end{document}